\backmatter
\begin{thebibliography}{99}
\bibitem{Hamilton1844} W.~R. Hamilton, On quaternions; or on a new system of imaginaries in algebra, \emph{Philosophical Magazine} \textbf{25} (1844) 10--13.
\bibitem{ConwaySmith} J.~H. Conway and D.~A. Smith, \emph{On Quaternions and Octonions} (A K Peters, 2003).
\bibitem{Hopf1931} H.~Hopf, \"Uber die Abbildungen der dreidimensionalen Sph\"are auf die Kugelfl\"ache, \emph{Math. Ann.} \textbf{104} (1931) 637--665.
\bibitem{Urbantke2003} H.~K. Urbantke, The Hopf fibration---seven times in physics, \emph{J. Geom. Phys.} \textbf{46} (2003) 125--150.
\bibitem{BengtssonZyczkowski} I.~Bengtsson and K.~\.Zyczkowski, \emph{Geometry of Quantum States}, 2nd ed. (Cambridge University Press, 2017).
\bibitem{GentiliStruppa2007} G.~Gentili and D.~C. Struppa, A new theory of regular functions of a quaternionic variable, \emph{Adv. Math.} \textbf{216} (2007) 279--301.
\bibitem{PenroseRindler1} R.~Penrose and W.~Rindler, \emph{Spinors and Space-Time}, Vol.~1 (Cambridge University Press, 1984).
\bibitem{Wigner1939} E.~P. Wigner, On unitary representations of the inhomogeneous Lorentz group, \emph{Ann. Math.} \textbf{40} (1939) 149--204.
\bibitem{Dirac1928} P.~A.~M. Dirac, The quantum theory of the electron, \emph{Proc. R. Soc. Lond. A} \textbf{117} (1928) 610--624.
\bibitem{Srednicki} M.~Srednicki, \emph{Quantum Field Theory} (Cambridge University Press, 2007).
\bibitem{Fueter1934} R.~Fueter, Die Funktionentheorie der Differentialgleichungen $\Delta u=0$ und $\Delta\Delta u=0$ mit vier reellen Variablen, \emph{Comment. Math. Helv.} \textbf{7} (1934) 307--330.
\bibitem{Sudbery1979} A.~Sudbery, Quaternionic analysis, \emph{Math. Proc. Camb. Phil. Soc.} \textbf{85} (1979) 199--225.
\bibitem{GHS} K.~G\"urlebeck, K.~Habetha and W.~Spr\"o\ss ig, \emph{Holomorphic Functions in the Plane and $n$-dimensional Space} (Birkh\"auser, 2008).
\bibitem{Lanczos1929} C.~Lanczos, Die tensoranalytischen Beziehungen der Diracschen Gleichung, \emph{Z. Phys.} \textbf{57} (1929) 447--473.
\bibitem{Davis} P.~J. Davis, \emph{Circulant Matrices} (Wiley, 1979).
\end{thebibliography}
